\begin{center}
    \chapter{Methodology of the Study}
\end{center}
\label{chap:methodology}

\noindent
The methodology chapter outlines the step-by-step process used to conduct
this research. It begins by describing how the data was prepared, including
dataset organization, missing-value handling, feature scaling, and temporal
splitting for training, validation, testing, and conformal calibration.
The chapter then explains the forecasting model architecture, focusing on
the Any-Quantile AQ-NBEATS backbone and our exogenous adaptive extension.
The core probabilistic forecasting framework follows the Any-Quantile
formulation proposed by Smyl et al.~\cite{smyl2025anyquantile}.
Key aspects such as quantile sampling strategies, stage-wise hyperparameter
configuration, and evaluation metrics are also presented. Overall, this
chapter provides a comprehensive understanding of how the study was
structured and which methodological choices were used to ensure reliable
and accurate probabilistic electricity demand forecasting.

\section{Methodological Framework}
\label{sec:method_framework}

Our methodology is organized as a reproducible research pipeline with five
sequential phases:

\begin{enumerate}
  \item \textbf{Data preparation}: acquisition, integrity checks,
        temporal splitting, and feature engineering.
  \item \textbf{Baseline reproduction}: replication of AQ-NBEATS settings
        from the reference study~\cite{smyl2025anyquantile}.
  \item \textbf{Model extension}: integration of exogenous features and
        quantile-sampling strategies in a staged protocol.
  \item \textbf{Uncertainty calibration}: conformal post-processing with
        CQR~\cite{romano2019cqr} to improve interval validity.
  \item \textbf{Evaluation and ablation}: metric-based comparison under
        identical data and training conditions.
\end{enumerate}

This structure ensures that each gain can be attributed to a specific
methodological change rather than to uncontrolled experimental factors.

\begin{figure}[H]
\centering
\includegraphics[width=0.95\textwidth]{methodology_final_reference_style.png}
\caption{Global methodology workflow used in this study.}
\label{fig:method_pipeline_placeholder}
\end{figure}
Figure~\ref{fig:method_pipeline_placeholder} presents a compact overview of
the full methodological pipeline and highlights the transition from
model-development stages to post-hoc conformal calibration.


\section{Data Collection and Understanding}
\label{sec:data_collection}

\subsection{Dataset Source}
\label{subsec:dataset_source}

The study uses the EMHIRES electricity demand dataset from the ENTSO-E
Transparency Platform (35 European countries, hourly granularity,
2006--2018). The resulting panel exhibits high heterogeneity in demand scale,
country-specific seasonality, and occasional missing values in late 2018 for
selected countries.

These properties make the task particularly suitable for quantile forecasting,
where both central tendency and uncertainty structure must be modeled across
distinct operating regimes.

\begin{figure}[H]
\centering
\includegraphics[width=0.95\textwidth]{dataset_characteristics_2018.png}
\caption{Illustrative heterogeneity of EMHIRES demand series across countries.}
\label{fig:data_heterogeneity_placeholder}
\end{figure}

Figure~\ref{fig:data_heterogeneity_placeholder} presents a multi-panel visualization highlighting the core challenges addressed by our modeling approach for the 2018 study period:

\subsection{Exogenous and Calendar Features}
\label{subsec:features}

To improve demand sensitivity to weather and temporal context, we augment the
target series with:

\begin{itemize}
  \item \textbf{Continuous meteorological covariates:} temperature,
        humidity, pressure, and wind speed.
  \item \textbf{Calendar covariates:} hour-of-day, day-of-week,
        month-of-year, and weekend indicator.
\end{itemize}

Calendar variables are deterministically known at forecast time and therefore
introduce no additional forecast uncertainty. Continuous weather features are
included to model short-term load sensitivity to atmospheric conditions.

In practice, this feature design allows the model to separate deterministic
calendar effects (e.g., weekday/weekend rhythm) from uncertain
weather-driven fluctuations, which is essential for stable probabilistic
intervals during peak-load and shoulder-load periods.

Formally, for each country-specific timestamp $t$, we define an exogenous
vector
\begin{equation}
\mathbf{e}_t = [\mathrm{temp}_t,\mathrm{hum}_t,\mathrm{press}_t,\mathrm{wind}_t,
\mathrm{hour}_t,\mathrm{dow}_t,\mathrm{month}_t,\mathrm{weekend}_t],
\label{eq:exog_vector}
\end{equation}
where continuous components are standardized and calendar components are
mapped to trainable embedding spaces before fusion.

\section{Dataset Organization and Splitting Protocol}
\label{sec:data_split}

The temporal split follows the reference protocol for fair comparison:

\begin{itemize}
  \item \textbf{Training:} 2006-01-01 to 2016-12-31.
  \item \textbf{Validation:} 2017-01-01 to 2017-12-31.
  \item \textbf{Test:} 2018-01-01 to 2018-12-31.
\end{itemize}

Each sample uses a lookback window of $W=168$ hours and a forecast horizon of
$H=48$ hours. Sliding windows are generated with a fixed stride to produce
supervised tuples of historical demand, historical/future exogenous features,
country identifier, and future targets.

For each sample index $s$, the supervised training tuple is:
\begin{equation}
\mathcal{D}_s=\left(\mathbf{x}_s,\mathbf{E}^{\mathrm{hist}}_s,\mathbf{E}^{\mathrm{fut}}_s,c_s,\mathbf{y}_s\right),
\label{eq:sample_tuple}
\end{equation}
with
\begin{equation}
\mathbf{x}_s=[y_{t-W+1},\dots,y_t],\quad
\mathbf{y}_s=[y_{t+1},\dots,y_{t+H}],
\label{eq:x_y_windows}
\end{equation}
where $c_s$ denotes the country index and $\mathbf{E}^{\mathrm{hist}}_s,
\mathbf{E}^{\mathrm{fut}}_s$ are the historical and future exogenous matrices.

For conformal calibration, the 2018 period is additionally partitioned into:

\begin{itemize}
  \item \textbf{Calibration split:} January-June 2018.
  \item \textbf{Conformal test split:} July-December 2018.
\end{itemize}

This split enforces strict non-leakage between calibration and final
interval evaluation.

\begin{figure}[H]
\centering
\includegraphics[width=0.95\textwidth]{temporal_splitting_protocol_fixed.png}
\caption{Chronological split design for training, validation, testing, and CQR calibration.}
\label{fig:split_protocol_placeholder}
\end{figure}

Figure~\ref{fig:split_protocol_placeholder} illustrates the strict temporal organization of data used throughout this study. The overall dataset spans 2006-2018 (13 years), and is divided into three non-overlapping periods to ensure an honest forward-in-time evaluation protocol with no leakage from future data:

\begin{itemize}
  \item \textbf{Training (2006--2016):} An 11-year window used for model fitting and learning of demand patterns, feature relationships, and seasonal dynamics.
  \item \textbf{Validation (2017):} A 1-year holdout period used for hyperparameter tuning and early stopping to select the best model checkpoint during training.
  \item \textbf{Test (2018):} A final 1-year holdout reserved for evaluation of forecasting performance on completely unseen data.
\end{itemize}

Within the test period (2018), an additional split is enforced for conformal calibration:

\begin{itemize}
  \item \textbf{CQR Calibration (January-June 2018):} Used to compute nonconformity scores and the empirical quantile threshold for interval correction (Eq.~\ref{eq:cqr_score_method}).
  \item \textbf{CQR Test (July-December 2018):} Held entirely separate and used only for the final evaluation of calibrated prediction intervals.
\end{itemize}

This chronological design enforces strict temporal consistency, all information used for training,

\section{Preprocessing Pipeline}
\label{sec:preprocess}

The preprocessing stage is fixed for all compared configurations:

\begin{enumerate}
  \item \textbf{Missing values:} training NaNs are imputed using weekly
        lag replacement ($t-168$) to preserve seasonality consistency.
  \item \textbf{Feature scaling:} continuous exogenous features are
        standardized with training statistics only.
  \item \textbf{Categorical encoding:} calendar indices are mapped through
        learnable embeddings.
  \item \textbf{Metric filtering:} NaN-marked timestamps in the test period
        are excluded from metric computation.
\end{enumerate}

The same preprocessing code path is used across all stages to guarantee
methodological comparability.

This unified preprocessing policy is critical for fair ablation studies,
because it ensures that observed gains originate from model or training
changes rather than from data handling differences.

\section{Modeling Strategy Based on Any-Quantile}
\label{sec:modeling_strategy}

\subsection{Baseline: AQ-NBEATS Reproduction}
\label{subsec:baseline_repro_method}

The baseline backbone is NBEATSAQCAT,
reproduced using the published hyperparameter family~\cite{smyl2025anyquantile}.
The model learns a direct mapping from context window and quantile level
$\tau \in (0,1)$ to horizon-wise quantile forecasts.

Let $f_{\theta}$ denote the forecaster. For one sample $s$ and one quantile
level $\tau$, the horizon output is:
\begin{equation}
\hat{\mathbf{q}}_{s,\tau}=f_{\theta}(\mathbf{x}_s,\mathbf{E}^{\mathrm{hist}}_s,
\mathbf{E}^{\mathrm{fut}}_s,c_s,\tau) \in \mathbb{R}^{H}.
\label{eq:forecast_mapping}
\end{equation}

\subsection{Extended Forecaster: Exogenous Adaptive Variant}
\label{subsec:exog_adaptive_model}

Our main extension uses
\texttt{AnyQuantileForecasterExogAdaptive}, which integrates:

\begin{itemize}
  \item historical and future exogenous features,
  \item calendar embeddings,
  \item adaptive quantile sampling during fine-tuning.
\end{itemize}

The objective remains quantile-consistent forecasting, with improved representation of demand drivers and better tail behavior under distributional shifts. Feature fusion can be summarized as:
\begin{equation}
\mathbf{z}_s = \phi_x(\mathbf{x}_s)+\phi_h(\mathbf{E}^{\mathrm{hist}}_s)
+\phi_f(\mathbf{E}^{\mathrm{fut}}_s)+\phi_c(c_s),
\label{eq:feature_fusion}
\end{equation}
where $\phi_x,\phi_h,\phi_f,\phi_c$ are learnable projections of demand history, historical exogenous variables, future exogenous variables, and country identity, respectively.

\begin{figure}[H]
\centering
\includegraphics[width=0.95\textwidth]{figure_4_4_architecture_en.png}
\caption{Conceptual architecture of the exogenous adaptive AQ-NBEATS forecaster.}
\label{fig:model_architecture_placeholder}
\end{figure}

Figure~\ref{fig:model_architecture_placeholder} illustrates how historical
load, exogenous variables, and calendar embeddings are fused and processed by
the AQ-NBEATS backbone to produce multi-quantile forecasts.

\subsection{Quantile Loss Formulation}
\label{subsec:quantile_loss}

For a target $y$ and predicted quantile $\hat{q}_{\tau}$, the pinball loss is:

\begin{equation}
\mathcal{L}_{\tau}(y, \hat{q}_{\tau}) =
(\tau - \mathbb{I}\{y < \hat{q}_{\tau}\})(y - \hat{q}_{\tau}).
\label{eq:pinball_method}
\end{equation}

Training minimizes an averaged multi-quantile objective over sampled levels
$\tau$ and forecast horizons. This is consistent with the Any-Quantile
paradigm and supports CRPS-oriented improvements.

Given a set of sampled quantiles $\mathcal{T}_s$ for sample $s$, the training
objective over a mini-batch $\mathcal{B}$ is:
\begin{equation}
\mathcal{L}_{\mathrm{MQ}}(\theta)=
\frac{1}{|\mathcal{B}|}
\sum_{s\in\mathcal{B}}
\frac{1}{|\mathcal{T}_s|H}
\sum_{\tau\in\mathcal{T}_s}\sum_{h=1}^{H}
\mathcal{L}_{\tau}(y_{s,h},\hat{q}_{s,h,\tau}).
\label{eq:mq_loss}
\end{equation}

Operationally, this objective enforces asymmetric penalties for under- and
over-estimation at each quantile level, enabling the model to learn complete
predictive distributions rather than only point forecasts.

\section{Stage-wise Training Protocol (Core Contribution)}
\label{sec:stage_protocol}

The proposed development follows a two-stage training curriculum,
then a calibration stage:

\subsection{Stage 1: Exogenous Foundation with Beta Sampling}
\label{subsec:stage1}

Stage 1 establishes a stable exogenous-aware foundation model using
AQNBEATSPlusPlus with q\_sampling=random\_in\_batch (beta-based random quantile sampling).
This stage aims to improve global distribution learning and tail exposure
while preserving optimization stability.

In this phase, quantile levels are sampled from a symmetric beta distribution:
\begin{equation}
\tau \sim \mathrm{Beta}(\beta,\beta),
\label{eq:beta_sampling}
\end{equation}
with $\beta<1$ to increase tail-sample frequency and improve lower/upper
quantile training.

\subsection{Stage 2: Adaptive Quantile Fine-Tuning}
\label{subsec:stage2}

Stage 2 initializes from Stage 1 weights
(loaded with \texttt{strict=False}) and switches to adaptive quantile
sampling. This training phase emphasizes poorly calibrated quantile regions,
improving sharpness and ranking quality. The best checkpoint is selected by
validation criterion (best epoch observed at epoch 4 in our final run).

If $g(\tau)$ denotes the current calibration gap estimate for quantile
$\tau$, adaptive sampling can be expressed as:
\begin{equation}
p(\tau)=\frac{\exp\left(\gamma\,|g(\tau)|\right)}
{\sum_{\tau'\in\mathcal{Q}}\exp\left(\gamma\,|g(\tau')|\right)},
\label{eq:adaptive_sampling}
\end{equation}
where $\gamma$ is a temperature parameter and $\mathcal{Q}$ is the candidate
quantile set. This mechanism allocates more updates to underperforming
quantile regions.


\subsection{Stage 3: Conformal Calibration (AQ-NBEATS-XA + CQR)}
\label{subsec:stage3_cqr}

To enforce reliable empirical coverage, we apply CQR on top of the
Stage-2 predictor. For lower/upper quantile predictions
$\hat{q}_{\alpha/2}(x_i)$ and $\hat{q}_{1-\alpha/2}(x_i)$, nonconformity scores
on calibration data are:

\begin{equation}
s_i = \max\left\{\hat{q}_{\alpha/2}(x_i)-y_i,\; y_i-\hat{q}_{1-\alpha/2}(x_i)\right\}.
\label{eq:cqr_score_method}
\end{equation}

Let $\hat{s}_{1-\alpha}$ be the empirical $(1-\alpha)$-quantile of
$\{s_i\}$. The calibrated interval is:

\begin{equation}
\left[\hat{q}_{\alpha/2}(x)-\hat{s}_{1-\alpha},\;
\hat{q}_{1-\alpha/2}(x)+\hat{s}_{1-\alpha}\right].
\label{eq:cqr_interval_method}
\end{equation}

This correction is model-agnostic and does not require retraining,
which makes it practical for deployment.

\begin{figure}[H]
\centering
\includegraphics[width=0.95\textwidth]{cqr_results.png}
\caption{CQR calibration results showing conformalized quantile intervals with improved coverage guarantees. The visualization demonstrates the effectiveness of post-hoc calibration in achieving nominal coverage levels across different quantile ranges.}
\label{fig:cqr_results}
\end{figure}

\section{Evaluation Protocol and Metrics}
\label{sec:eval_protocol}

All configurations are evaluated under identical splits and preprocessing.
Let $N$ denote the number of evaluated observations and $\mathcal{T}$ the
set of quantile levels used for scoring.

\subsection{CRPS and Quantile Approximation}
\label{subsec:crps_formal}

CRPS is the primary ranking metric for probabilistic forecasts~\cite{gneiting2007strictly}.
Using quantile forecasts, it is approximated by:
\begin{equation}
\mathrm{CRPS} \approx \frac{2}{|\mathcal{T}|N}
\sum_{i=1}^{N}\sum_{\tau\in\mathcal{T}}
\mathcal{L}_{\tau}(y_i,\hat{q}_{i,\tau}).
\label{eq:crps_quantile}
\end{equation}

\subsection{Normalized CRPS}
\label{subsec:ncrps_formal}

To compare across countries with very different demand scales, we use:
\begin{equation}
\mathrm{N\text{-}CRPS}=\frac{\mathrm{CRPS}}{\bar{y}},
\label{eq:ncrps_formal}
\end{equation}
where $\bar{y}$ is the mean load over the corresponding evaluation segment.

\subsection{Coverage of Prediction Intervals}
\label{subsec:coverage_formal}

For a target miscoverage level $\alpha$ (e.g., $\alpha=0.05$ for 95\% PI),
empirical coverage is computed as:
\begin{equation}
\mathrm{Coverage}_{1-\alpha}=\frac{1}{N}\sum_{i=1}^{N}
\mathbb{I}\left\{\hat{q}_{i,\alpha/2}\le y_i\le \hat{q}_{i,1-\alpha/2}\right\}.
\label{eq:coverage_formal}
\end{equation}

\subsection{Point Forecast Metrics}
\label{subsec:point_metrics_formal}

Point metrics are computed from the median forecast
$\hat{y}_i=\hat{q}_{i,0.5}$:
\begin{align}
\mathrm{MAE} &= \frac{1}{N}\sum_{i=1}^{N}|y_i-\hat{y}_i|, \label{eq:mae_formal} \\
\mathrm{MAPE} &= \frac{100}{N}\sum_{i=1}^{N}\left|\frac{y_i-\hat{y}_i}{y_i}\right|, \label{eq:mape_formal} \\
\mathrm{SMAPE} &= \frac{100}{N}\sum_{i=1}^{N}
\frac{2|y_i-\hat{y}_i|}{|y_i|+|\hat{y}_i|}. \label{eq:smape_formal}
\end{align}

For robustness, MAPE is computed on non-zero targets only.

\subsection{Reporting Protocol}
\label{subsec:reporting_protocol}

For each configuration, we report the primary CRPS ranking value and
complementary calibration and point metrics under the same filtered test set
(NaN timestamps excluded).

To reduce randomness bias, we track checkpointed validation performance and
report final values on the untouched test period. Calibration and test sets
are strictly separated for CQR.

For publication reporting, each metric should be accompanied by the exact
evaluation split definition and, whenever available, variability indicators
(standard deviation or confidence intervals) to support robust interpretation.

\section{Methodological Traceability and Reproducibility}
\label{sec:reproducibility}

Reproducibility is ensured through:

\begin{itemize}
  \item fixed data split definitions,
  \item version-controlled configuration files and code,
  \item deterministic preprocessing and feature pipelines,
  \item explicit stage-to-stage checkpoint transfer,
  \item standardized metric scripts shared across experiments.
\end{itemize}

The methodology is therefore auditable from raw data ingestion to final
probabilistic outputs.

\section{Methodology Outcome Summary}
\label{sec:method_outcomes}

This section provides a concise summary of the developed methodological
pipeline and the corresponding performance outcomes, with a specific focus
on the Continuous Ranked Probability Score (CRPS) achieved at each major
stage. The sequential application of our main contributions---exogenous
feature integration with beta sampling, adaptive quantile fine-tuning, and
post-hoc Conformal Quantile Regression (CQR)---defines a structured pathway
for improving probabilistic forecasting accuracy.

Table~\ref{tab:method_stage_summary} (Table~4.1) summarizes the final
methodological stages and their respective CRPS values, highlighting the
incremental gains achieved throughout the project.

\begin{table}[H]
\centering
\caption{Stage-wise methodology and obtained CRPS values.}
\label{tab:method_stage_summary}
\begin{tabular}{llc}
\toprule
\textbf{Stage} & \textbf{Configuration} & \textbf{CRPS} \\
\midrule
Stage 1 (Final) & Exog + Beta sampling foundation & 178.16 \\
Stage 2 (Final) & Exog + Adaptive sampling fine-tuning & 169.39 \\
AQNBEATS++ + CQR & Conformalized calibrated version & \textbf{126.66} \\
\bottomrule
\end{tabular}
\end{table}

The results presented in Table~\ref{tab:method_stage_summary} demonstrate a
consistent and substantial improvement over baseline behavior. Stage~1,
which establishes an exogenous-aware foundation model using beta sampling,
reduces CRPS to 178.16~MW. This gain confirms the value of stronger global
distribution learning and improved tail exposure. Stage~2 further refines
the model through adaptive quantile sampling and reaches 169.39~MW,
indicating improved sharpness and better correction of poorly calibrated
quantile regions.

The final AQNBEATS++ + CQR stage delivers the largest performance jump,
achieving a CRPS of 126.66~MW. This result confirms the effectiveness of
post-hoc conformal calibration for enforcing reliable empirical coverage
while further improving forecast interval quality. Overall, the full
methodological sequence---robust exogenous modeling, adaptive quantile
refinement, and conformal calibration---proves highly effective for
probabilistic short-term load forecasting and establishes state-of-the-art
performance in this study setting.

